% llm_embeddings_lexsem.tex
% Compile: pdflatex llm_embeddings_lexsem.tex  (or lualatex/xelatex)
% If you want bibliography: run biber (see notes near the end)

\documentclass{beamer}
\input{shared.tex}

\usepackage{booktabs}


% --- Optional: biblatex (uncomment if you want refs printed) ---
% \usepackage[backend=biber,style=authoryear]{biblatex}
% \addbibresource{llm_embeddings_lexsem.bib}

\title[Lexical meaning in the LLM era]{Word Meaning, Embeddings, and LLMs\\(with BERT/ELMo and Multilingual Embeddings)}
\author{Francis Bond}
\institute{Methods in Lexical Semantics (LOT / RM / PhD)}
\date{}

\begin{document}

% 1
\begin{frame}
  \titlepage
\end{frame}

% 2
\begin{frame}{What this lecture is about}
\begin{itemize}
  \item How do modern systems represent word meaning?
  \item From symbolic lexical resources to vectors to contextual LLM representations
  \item What you can \emph{measure} (and what you should \emph{distrust})
\end{itemize}
\end{frame}

% 3
\begin{frame}{Learning outcomes}
By the end, you should be able to:
\begin{itemize}
  \item explain the difference between \textbf{static} and \textbf{contextual} embeddings
  \item connect embedding behavior to classic lexical semantic phenomena (polysemy, similarity, relations)
  \item outline how multilingual embeddings enable cross-lingual comparison
  \item design a small evaluation or probing task using embeddings
\end{itemize}
\end{frame}

% 4
\begin{frame}{Plan (high level)}
\begin{enumerate}
  \item Lexical meaning: what we want to capture
  \item Static embeddings: distributional semantics
  \item Contextual embeddings: ELMo $\rightarrow$ BERT $\rightarrow$ LLMs
  \item Multilingual embeddings: alignment + cross-lingual transfer
  \item WordNets vs embeddings: complements, hybrids, and pitfalls
\end{enumerate}
\end{frame}

% 5
\begin{frame}{Lexical semantics refresher: the targets}
\begin{itemize}
  \item \textbf{Sense} vs \textbf{use}: lexeme, lemma, token
  \item \textbf{Polysemy}: systematic vs accidental
  \item \textbf{Lexical relations}: synonymy, hypernymy, meronymy, antonymy
  \item \textbf{Context effects}: meaning modulation, coercion, selectional preferences
\end{itemize}
\end{frame}

% 6
\begin{frame}{Why representations matter}
\begin{itemize}
  \item We use representations to:
  \begin{itemize}
    \item test hypotheses (e.g., semantic neighborhoods, category structure)
    \item generalize (e.g., inference, similarity, analogy)
    \item build systems (WSD, IR, MT, QA, IE)
  \end{itemize}
  \item Different representations make different assumptions visible/invisible
\end{itemize}
\end{frame}

% 7
\begin{frame}{Symbolic vs distributional: a mental contrast}
\begin{itemize}
  \item \textbf{Symbolic} (WordNet): explicit senses and relations, interpretable edges
  \item \textbf{Distributional} (embeddings): similarity from usage patterns, implicit structure
  \item \textbf{Contextual} (BERT/LLMs): meaning as a function of context, dynamic vectors
\end{itemize}
\end{frame}

% 8
\begin{frame}{The distributional hypothesis (intuition)}
\begin{itemize}
  \item ``You shall know a word by the company it keeps''
  \item Words with similar contexts tend to have related meanings
  \item Build meaning from co-occurrence statistics
\end{itemize}
\end{frame}

% 9
\begin{frame}{From counts to vectors (very schematic)}
\begin{itemize}
  \item Choose a vocabulary of \textbf{targets} and \textbf{contexts}
  \item Fill a co-occurrence matrix (counts or weighted)
  \item Reduce dimensionality / learn parameters
  \item Use cosine similarity to compare vectors
\end{itemize}
\end{frame}

% 10
\begin{frame}{Static embeddings: what they are}
\begin{itemize}
  \item One vector per word type (e.g., \texttt{bank} has one vector)
  \item Learned to predict contexts (or from matrix factorization)
  \item Captures graded similarity and topical association
  \item Struggles with polysemy and compositional nuance
\end{itemize}
\end{frame}

% 11
\begin{frame}{Word2vec: CBOW and Skip-gram (idea)}
\begin{itemize}
  \item CBOW: predict target from surrounding words
  \item Skip-gram: predict surrounding words from target
  \item Training yields embeddings that capture distributional regularities
\end{itemize}
\end{frame}

% 12
\begin{frame}{What ``similarity'' in embeddings often reflects}
\begin{itemize}
  \item taxonomic similarity (dog $\sim$ cat)
  \item topical association (dog $\sim$ leash)
  \item register/domain clustering
  \item frequency effects (high-frequency words become hubs)
\end{itemize}
\end{frame}

% 13
\begin{frame}{Mini-exercise (discussion)}
Pick a word: \texttt{paper}, \texttt{mouse}, \texttt{bank}.
\begin{itemize}
  \item What different senses/usages do you expect?
  \item In a static embedding, which neighbors would you expect?
  \item Would they group by sense, topic, or frequency?
\end{itemize}
\end{frame}

% 14
\begin{frame}{Evaluating static embeddings (classic)}
\begin{itemize}
  \item word similarity datasets (human judgments)
  \item analogy tasks (limited diagnostic value)
  \item downstream tasks (task-specific, can hide problems)
\end{itemize}
\end{frame}

% 15
\begin{frame}{Why lexical semantics cares about failure cases}
\begin{itemize}
  \item Polysemy: one vector conflates senses
  \item Antonymy: distributionally similar but semantically opposed
  \item Rare senses: dominated by majority usage
  \item Multiword expressions: meanings not reducible to parts
\end{itemize}
\end{frame}

% 16
\begin{frame}{From static to contextual: the key shift}
\begin{itemize}
  \item Static: vector(word)
  \item Contextual: vector(word, context)
  \item Same surface word gets different embeddings in different sentences
\end{itemize}
\end{frame}

% 17
\begin{frame}{ELMo (high-level)}
\begin{itemize}
  \item Deep contextualized word representations from a language model
  \item BiLSTM: left and right contexts
  \item Representations are layer-wise; different layers encode different info
\end{itemize}
\end{frame}

% 18
\begin{frame}{ELMo: why it mattered}
\begin{itemize}
  \item Demonstrated strong gains on many NLP tasks
  \item Provided an operational way to represent polysemy via context
  \item Encouraged probing analyses: ``what is encoded where?''
\end{itemize}
\end{frame}

% 19
\begin{frame}{Transformers: the next leap}
\begin{itemize}
  \item Self-attention replaces recurrent sequence processing
  \item Better long-range dependencies
  \item Efficient parallel training at scale
\end{itemize}
\end{frame}

% 20
\begin{frame}{BERT (high-level)}
\begin{itemize}
  \item Bidirectional Transformer encoder
  \item Pretraining objectives:
  \begin{itemize}
    \item masked language modeling (MLM)
    \item (originally) next sentence prediction (NSP)
  \end{itemize}
  \item Produces contextual token embeddings
\end{itemize}
\end{frame}

% 21
\begin{frame}{Tokenization matters (WordPiece/BPE)}
\begin{itemize}
  \item Many models operate on subword units
  \item ``word meaning'' becomes distributed over tokens
  \item Morphology interacts with representation (good and bad)
\end{itemize}
\end{frame}

% 22
\begin{frame}{Contextual embeddings and word senses}
\begin{itemize}
  \item Clustering token embeddings of the same word often yields sense-like groupings
  \item But clusters depend on:
  \begin{itemize}
    \item corpus domain and genre
    \item model architecture and layer
    \item tokenization and frequency
  \end{itemize}
\end{itemize}
\end{frame}

% 23
\begin{frame}{Which layer should I use?}
\begin{itemize}
  \item Lower layers: morphology, local syntax, surface patterns
  \item Middle layers: syntax/semantics interface signals
  \item Higher layers: task-like abstractions, discourse-ish effects
  \item Rule of thumb: try several layers and report sensitivity
\end{itemize}
\end{frame}

% 24
\begin{frame}{Probing lexical semantic information in BERT-like models}
\begin{itemize}
  \item similarity neighborhoods for contextual tokens
  \item relation classification (hypernymy/meronymy probing)
  \item selectional preference tests (verb-object plausibility)
  \item sense discrimination via clustering / contrastive contexts
\end{itemize}
\end{frame}

% 25
\begin{frame}{LLMs: where do embeddings fit?}
\begin{itemize}
  \item LLMs still begin with an \textbf{embedding layer} for tokens
  \item Hidden states evolve through many layers (contextual representations)
  \item Output layer maps back to vocabulary logits (or tied embeddings)
\end{itemize}
\end{frame}

% 26
\begin{frame}{Word meaning from the LLM viewpoint}
\begin{itemize}
  \item Meaning as \textbf{predictive utility}: what helps next-token prediction?
  \item Representations are:
  \begin{itemize}
    \item highly contextual
    \item multi-purpose (syntax, semantics, world knowledge signals)
    \item shaped by training data distributions and objectives
  \end{itemize}
\end{itemize}
\end{frame}

% 27
\begin{frame}{Lexical semantics in an LLM pipeline}
\begin{itemize}
  \item Retrieval + embeddings (semantic search)
  \item Prompting + in-context learning
  \item Fine-tuning / adapters
  \item Evaluation: semantic consistency across paraphrases and languages
\end{itemize}
\end{frame}

% 28
\begin{frame}{Static vs contextual: a quick comparison table}
\begin{tabular}{@{}lll@{}}
\toprule
 & Static embeddings & Contextual embeddings \\
\midrule
Representation & word type & token-in-context \\
Polysemy & conflated & partially separated \\
Interpretability & low & low (but probeable) \\
Compute & cheap & expensive \\
Cross-lingual & alignable & depends on model \\
\bottomrule
\end{tabular}
\end{frame}

% 29
\begin{frame}{Multilingual embeddings: why linguists care}
\begin{itemize}
  \item compare lexicalization patterns across languages
  \item transfer models/data (low-resource settings)
  \item measure semantic shift and contact effects cross-lingually
  \item support multilingual resources (WordNets, dictionaries)
\end{itemize}
\end{frame}

% 30
\begin{frame}{Two main routes to multilingual embeddings}
\begin{enumerate}
  \item \textbf{Align monolingual spaces} (post-hoc mapping)
  \item \textbf{Train multilingual models} (shared parameters/data)
\end{enumerate}
\end{frame}

% 31
\begin{frame}{Aligned monolingual embeddings (sketch)}
\begin{itemize}
  \item Train embeddings separately per language
  \item Learn a mapping using:
  \begin{itemize}
    \item bilingual lexicon / seed dictionary
    \item identical strings / numerals
    \item unsupervised alignment assumptions
  \end{itemize}
  \item Result: shared vector space for cross-lingual similarity
\end{itemize}
\end{frame}

% 32
\begin{frame}{FastText: a practical ingredient}
\begin{itemize}
  \item Subword-based embeddings help morphology and rare words
  \item Commonly used as a baseline for multilingual alignment
  \item Easy to get for many languages
\end{itemize}
\end{frame}

% 33
\begin{frame}{MUSE-style alignment (intuition)}
\begin{itemize}
  \item Find a linear map between spaces
  \item Iteratively refine with nearest-neighbor lexicons
  \item Evaluate with bilingual lexicon induction
\end{itemize}
\end{frame}

% 34
\begin{frame}{Multilingual contextual models}
\begin{itemize}
  \item mBERT: multilingual BERT trained on many Wikipedias
  \item XLM-R: robust multilingual pretraining on large web data
  \item LaBSE / LASER: sentence-level multilingual representations
\end{itemize}
\end{frame}

% 35
\begin{frame}{What ``shared space'' means in mBERT/XLM-R}
\begin{itemize}
  \item One model, shared subword vocabulary (partly)
  \item Languages can align implicitly through shared parameters
  \item Degree of alignment varies by:
  \begin{itemize}
    \item script overlap, data size, typological distance, domain
  \end{itemize}
\end{itemize}
\end{frame}

% 36
\begin{frame}{Cross-lingual lexical semantics with multilingual models}
\begin{itemize}
  \item nearest-neighbor translation (token or sentence embeddings)
  \item sense discrimination across languages (clustering)
  \item lexical gap detection (no close counterpart)
  \item semantic domain comparison (field-level analysis)
\end{itemize}
\end{frame}

% 37
\begin{frame}{Mini-exercise idea (multilingual embeddings)}
\begin{itemize}
  \item Pick 10 synsets from OMW (e.g., emotions, motion verbs)
  \item For each language, collect lemmas
  \item Compare:
  \begin{itemize}
    \item WordNet graph distance vs embedding cosine similarity
    \item Nearest neighbors in each language space
  \end{itemize}
\end{itemize}
\end{frame}

% 38
\begin{frame}{Where multilingual embeddings go wrong (often)}
\begin{itemize}
  \item language-specific polysemy patterns misalign
  \item named entities and cultural terms cluster oddly
  \item domain mismatch (Wikipedia vs conversational text)
  \item scripts and tokenization (segmentation differences)
\end{itemize}
\end{frame}

% 39
\begin{frame}{WordNet: why lexical graphs still matter}
\begin{itemize}
  \item explicit relations: hypernymy, meronymy, antonymy, etc.
  \item interpretable structure: paths, subgraphs, constraints
  \item supports controlled generalization and explainability
  \item useful for evaluation and error analysis of embeddings/LLMs
\end{itemize}
\end{frame}

% 40
\begin{frame}{WordNet vs embeddings: complementary strengths}
\begin{itemize}
  \item WordNet excels at \textbf{typed relations} and \textbf{sense inventories}
  \item Embeddings excel at \textbf{graded similarity} and \textbf{coverage} (if data exists)
  \item Together:
  \begin{itemize}
    \item sense clustering + labeling
    \item relation discovery + verification
    \item cross-lingual mapping sanity checks
  \end{itemize}
\end{itemize}
\end{frame}

% 41
\begin{frame}{A simple hybrid picture}
\begin{itemize}
  \item Use embeddings to propose candidates:
  \begin{itemize}
    \item synonyms / hypernyms / translations
  \end{itemize}
  \item Use WordNet/OMW to:
  \begin{itemize}
    \item validate relation type
    \item control sense granularity
    \item provide glosses and structure for interpretation
  \end{itemize}
\end{itemize}
\end{frame}

% 42
\begin{frame}{Embedding similarity vs WordNet distance}
\begin{itemize}
  \item WordNet distance: graph shortest path / information content
  \item Embedding distance: cosine similarity in vector space
  \item They correlate sometimes, but diverge:
  \begin{itemize}
    \item topical association vs taxonomic similarity
    \item antonymy and co-hyponymy
  \end{itemize}
\end{itemize}
\end{frame}

% 43
\begin{frame}{Polysemy: a worked example (template)}
\begin{itemize}
  \item Word: \texttt{bank}
  \item WordNet senses:
  \begin{itemize}
    \item financial institution
    \item river bank
  \end{itemize}
  \item Compare:
  \begin{itemize}
    \item static neighbors
    \item BERT token neighbors in two contexts
  \end{itemize}
\end{itemize}
\end{frame}

% 44
\begin{frame}{Contextual neighbors: how to do it (conceptually)}
\begin{itemize}
  \item Choose two sentences with different senses
  \item Extract token embedding from a chosen layer
  \item Retrieve nearest neighbors (in a reference bank of token embeddings)
  \item Interpret: does it separate sense or topic?
\end{itemize}
\end{frame}

% 45
\begin{frame}{Multilingual sense comparison (template)}
\begin{itemize}
  \item Pick a synset in OMW (concept anchor)
  \item List lemmas across languages
  \item Ask:
  \begin{itemize}
    \item Are lemmas equally polysemous?
    \item Are there lexical gaps?
    \item Does embedding similarity preserve the mapping?
  \end{itemize}
\end{itemize}
\end{frame}

% 46
\begin{frame}{LLMs and lexical relations: promise and risk}
\begin{itemize}
  \item LLMs can generate plausible relations (hypernyms, synonyms)
  \item But:
  \begin{itemize}
    \item relation \emph{type} can drift
    \item sense boundaries can blur
    \item hallucinations are confident and well-formed
  \end{itemize}
\end{itemize}
\end{frame}

% 47
\begin{frame}{A practical workflow for students}
\begin{enumerate}
  \item Use OMW as the semantic backbone (concepts + relations)
  \item Use embeddings to explore neighborhoods and candidate links
  \item Validate by:
  \begin{itemize}
    \item gloss inspection and examples
    \item relation constraints (POS, domain, hierarchy)
    \item cross-lingual consistency checks
  \end{itemize}
\end{enumerate}
\end{frame}

% 48
\begin{frame}{Multilingual embeddings: evaluation ideas}
\begin{itemize}
  \item bilingual lexicon induction accuracy
  \item cross-lingual semantic similarity judgment
  \item zero-shot transfer on a semantic task
  \item alignment stability across domains (news vs wiki vs social)
\end{itemize}
\end{frame}

% 49
\begin{frame}{What to report in a small study}
\begin{itemize}
  \item data source and preprocessing (tokenization matters!)
  \item model choice (static vs contextual; which multilingual model)
  \item layer choice and sensitivity analysis
  \item qualitative error analysis (with examples)
\end{itemize}
\end{frame}

% 50
\begin{frame}{Common pitfalls (lexical semantics edition)}
\begin{itemize}
  \item confusing topical association with synonymy
  \item treating embedding clusters as ``true senses''
  \item ignoring frequency and domain imbalance
  \item evaluating only on English-like phenomena
\end{itemize}
\end{frame}

% 51
\begin{frame}{Ethics and bias (brief)}
\begin{itemize}
  \item embeddings inherit biases in training data
  \item multilingual models can amplify majority-language norms
  \item lexical resources also encode editorial choices and cultural perspectives
\end{itemize}
\end{frame}

% 52
\begin{frame}{Where WordNet helps with bias analysis}
\begin{itemize}
  \item explicit semantic categories for auditing
  \item controlled contrasts across relation types
  \item sense-level framing can reduce some confounds (sometimes)
\end{itemize}
\end{frame}

% 53
\begin{frame}{Where embeddings help WordNet work}
\begin{itemize}
  \item suggest missing lemmas and candidate senses
  \item detect inconsistent placements (outliers in semantic fields)
  \item support semi-automatic linking across languages
\end{itemize}
\end{frame}

% 54
\begin{frame}{One slide on ``LLM-era'' takeaways}
\begin{itemize}
  \item Word meaning is now routinely modeled as \textbf{contextual representation}
  \item Lexical semantics remains essential for:
  \begin{itemize}
    \item evaluation, interpretability, and controlled inference
    \item cross-lingual conceptual comparison
    \item building and maintaining lexical resources
  \end{itemize}
\end{itemize}
\end{frame}

% 55
\begin{frame}{Suggested practical demo (in class)}
\begin{itemize}
  \item Choose 1 ambiguous word (EN + NL)
  \item Two contexts per sense
  \item Compare:
  \begin{itemize}
    \item OMW synsets + relations
    \item static neighbors (fastText)
    \item contextual neighbors (mBERT/XLM-R token embeddings)
  \end{itemize}
\end{itemize}
\end{frame}

% 56
\begin{frame}{Suggested mini-project prompt (1 page)}
\begin{itemize}
  \item Build a 10–15 concept mini-field via OMW
  \item Provide:
  \begin{itemize}
    \item a small graph visualization
    \item 2–3 cross-lingual ``interesting cases''
    \item embedding-based comparison + error analysis
  \end{itemize}
\end{itemize}
\end{frame}

% 57
\begin{frame}{If you only remember 5 points}
\begin{enumerate}
  \item Static embeddings conflate senses
  \item Contextual embeddings partially separate senses
  \item Tokenization and layer choice matter
  \item Multilingual spaces are imperfect and uneven
  \item WordNets provide interpretable structure for testing hypotheses
\end{enumerate}
\end{frame}

% 58
\begin{frame}{Further reading (core)}
\begin{itemize}
  \item ELMo (contextualized embeddings)
  \item BERT (transformer contextual representations)
  \item Multilingual models: mBERT, XLM-R, LASER/LaBSE
  \item WordNet/OMW docs for symbolic grounding
\end{itemize}
\end{frame}

% 59
\begin{frame}{References (optional slide)}
\begin{itemize}
  \item Add your preferred references here (or enable biblatex at top)
  \item I included a ready-to-use .bib below this file (copy/paste)
\end{itemize}
\end{frame}

% 60
\begin{frame}{Questions}
\centering
\Large Questions / discussion
\end{frame}

% --- If using biblatex, uncomment:
% \begin{frame}[allowframebreaks]{References}
%   \printbibliography
% \end{frame}

\end{document}

%%% Local Variables: 
%%% coding: utf-8
%%% mode: latex
%%% TeX-PDF-mode: t
%%% TeX-engine: lualatex
%%% End: 
